
\documentclass[../../../../main.tex]{subfiles}
% \documentclass[a4paper,10pt]{ltjsarticle}
\begin{document}
% \documentclass[a4paper,12pt]{article}
% \usepackage[utf8]{inputenc}
% \usepackage[japanese]{babel}
% \usepackage{xeCJK}
% \usepackage{geometry}
% \usepackage{array}
% \usepackage{booktabs}
% \usepackage{fancybox}

% % ページ設定
% \geometry{top=25mm, bottom=25mm, left=20mm, right=20mm}
% \pagestyle{empty}

% % フォント設定
% \setCJKmainfont{Noto Sans CJK JP}

% \begin{document}

% 年度の大きな表示
\begin{center}
  {\fontsize{48pt}{58pt}\selectfont \textbf{2010}\normalsize\textbf{年度}}
\end{center}

\vspace{10mm}

% 本文


\vspace{15mm}

% 出題テーマと難易度の表
\noindent\textbf{出題テーマと難易度}

\vspace{5mm}

\begin{center}
  \shadowbox{%
    \begin{tabular}{|c|p{8cm}|c|}
      \hline
      \textbf{問題番号} & \textbf{テーマ} & \textbf{難易度} \\
      \hline
      第1問           & 数列，極限        & 易            \\
      \hline
      第2問           & 平面図形         & 普            \\
      \hline
    \end{tabular}%
  }
\end{center}

\vspace{15mm}

% 解答
\noindent\textbf{解答}

\vspace{5mm}

\begin{center}
  \shadowbox{%
    \renewcommand{\arraystretch}{1.6}
    \begin{tabular}{|c|c|p{8cm}|}
      \hline
      \textbf{問題番号}        & \textbf{小問} & \textbf{答え}                                                                                                              \\
      \hline
      \multirow{2}{*}{第1問} & (1)         & $\displaystyle a_n = \frac{1}{2}(a+b)t^{n-1} + \frac{1}{2}(a-b)\left(\frac{10}{t^2+1}\right)^{n-1}$                      \\
      \cline{2-3}
                           & (2)         & $\displaystyle (a+b=0 \land (t \le -3 \text{ or } 3 \le t)) \text{ または } (a=0 \land t = 2)$                              \\                                                                                                  \\
      \hline
      \multirow{3}{*}{第2問} & (1)         & グラフを描く問題                                                                                                                 \\
      \cline{2-3}
                           & (2)         & $\displaystyle     n(\alpha) =
        \begin{dcases}
          0 < \alpha \le 1, \alpha = e & \cdots \text{1つ} \\
          1 < \alpha (\alpha \ne e)    & \cdots \text{2つ}
        \end{dcases}$
      \\
      \cline{2-3}
                           & (3)         & $\displaystyle S(\alpha)= 2 -\alpha (\log \alpha)^2 + 2\alpha \log \alpha - 2\alpha + \frac{1}{4}\alpha (\log \alpha)^3$ \\
      \hline
    \end{tabular}%
    \renewcommand{\arraystretch}{1.0}
  }
\end{center}

% 解答時間
\noindent\textbf{試験形態}

\vspace{5mm}

\begin{center}
  \shadowbox{%
    \begin{tabular}{p{8cm}}
      2問90分
    \end{tabular}%
  }
\end{center}


\end{document}